% Transcription — fonds Alexandre Grothendieck, shelfmark 58, batch 4 (pages 61-66).
% Established from the facsimile archives/batches/58/batch-04.pdf (PDF page k =
% archivists' page 60+k, no cover sheet), cut from a copy of 58.pdf fetched
% through the project's own relay
% (https://grothendieck-relay.onrender.com/source/58.pdf); tiles in
% archives/tiles/58/p61-p66, re-cropped at 350-400 dpi with pdftoppm. Per the
% skill transcribe-grothendieck. The folder is sixty-six pages in four batches
% (1-20, 21-40, 41-60, 61-66), transcribed in parallel; this is the last, a
% short batch of six pages, written after batch 2 (whose notation it keeps)
% and without sight of batch 3.
%
% Pass: Opus 5.5 (claude-opus-5-5), 2026-09-26 — first pass, unchecked against the pages by a human.
%
% Pages skipped: 62 — an upside-down typed leaf (page « - 7 - » of a paper on
% characters of finite groups: condition (B_q), Proposition 2 after Brauer
% and Witt, formulas (12)-(13), « 3. Démonstration du théorème 3 »), with
% nothing of his but a typist's insertion caret; skipped by the settled rule
% for typed versos. 63 — a cream sheet carrying only the offset of page 60
% (typescript and pencil figures transferred by contact), nothing written on
% it.
%
% Pages summarised: none. Pages transcribed: 61 (blue ink over the faint
% show-through of a typescript), 64 and 65 (two sheets of scratch work in blue
% and black ink: diagrams, formulas and small figures, with little prose),
% 66 (blue ink, a fair start of an argument breaking off mid-sentence).
%
% His pagination: none visible. His numbered runs: none; page 61 labels two
% groups of conditions « C'_i » and « C_n » with braces; page 66 cites
% « II 8.7.5 » and « II 8.8.6 » (presumably EGA II, § 8 on cones; the
% reference is his and has not been checked).
%
% What the batch is about. P. 61: for U = X - Z and g : U -> X, a chain of
% equivalences between the vanishing of H^i(U, F(m)) for m large, of the
% local cohomology sheaves H^{i+1}_Z(F) (i.e. H^{i+1}_z(F_z) = 0 at z in Z)
% and of R^i g_*(F), grouped as conditions C'_i; then C_n, vanishing for
% 0 < i <= n, equivalent (for F = g_*(F|U)) to prof_z F > n+1 at the points
% of Z; with the remark that only the conditions C require anything at the
% points of Z = X - U and not only at the points of U. P. 64: scratch on
% Pic of the projective cone C^ = Spec(?) S[z] over X, O(1), the blown-up
% cone C'^, the sequences Pic(X) -> Pic(C'^) -> Pic(C') -> Pic(C'~) with
% injectivity/surjectivity « dans le cas X loc. factoriel », Pic(X)/Z xi in
% Pic(C'^)/Z xi, a boxed question « Pic(C) = 0 ?? », and on the right the
% relations g^m = u f^n, g = u f^k. P. 65: the local cohomology sequence
% of O_X^* for U = X - a (0 -> H^0_a -> H^0(X) -> H^0(U) -> H^1_a -> H^1(X)
% -> H^1(U) -> H^2_a -> H^2(X)); the implications loc. factoriel =>
% loc. quasi-factoriel => parafactoriel, factoriel => 1-factoriel =>
% 0-factoriel; genus formulas for a curve with nodes (1 - g = sum(1 - g_i)
% - sum n_s, g = 1 + sum g_i + sum n_s - nu, the cases g = g' + sum n_s
% and g = (g' + g'') + sum n_s - 1) with figures of nodal crossings and
% k[x,y,z]/xy. P. 66: X projective over k with Gamma(X, O_X) = k, C the
% affine projecting cone, C^ the projective one; X --i--> X' --pi--> X,
% f : X' -> C^; Pic(X') = Pic(X) x Z, the kernel of i^* is Z generated by
% O_{X'}(1), Pic(C^) -> Pic(X') injective since f_*(O_{X'}) = O_{C^}
% « par II 8.8.6 », boxed Pic(C^) -> Ker[Pic(X') -> Pic(X)]; then C_X, the
% part of X' over C, and the claim « Je dis que : Pic(X) -> Pic(C_X) (?) »,
% with a margin « faut-il supposer X localement factoriel ?? »; the
% injectivity is argued and the surjectivity begun, breaking off.
%
% Notation, kept from batch 2 where it recurs: \underline{\mathcal{O}} for
% his underlined O (structure sheaves), \underline{H} for the sheaf of local
% cohomology; \hat{C} for the projective cone (his C with a hat), C' and
% \hat{C}' for the primed cones of p. 64, \widetilde{C'} for his C' with a
% tilde; \mathbf{Z} for his double-struck or bold Z (the integers), Z plain
% for the closed set X - U of p. 61; \xi for the letter after Z on p. 64.
%
% Readings to check first: the prose of p. 61 (« la profondeur », « quelque
% chose »), the bracket « [\ill{} supposant que F = g_*(F|U)] » and the
% overwritten arrow above it; « Spec » on p. 64 (Proj would be expected for a
% projective cone); the labels injectif/surjectif around the diagrams of
% p. 64 and which arrow each belongs to; the struck and overwritten terms of
% the long exact sequence of p. 65; the first term of « 1 - h' = 2 - h' -
% h'' » on p. 65; the whole prose of p. 66 below the box.

\documentclass[11pt,a4paper]{article}
% Le préambule commun à toutes les transcriptions du fonds Grothendieck.
%
% Il tient en une idée : **l'incertitude se compose, elle ne se résout pas.**
% Une transcription qui lisse un mot douteux a détruit la seule chose pour
% laquelle on la faisait — un lecteur doit pouvoir savoir, à chaque endroit,
% ce qui était sur le feuillet et ce qui a été ajouté. Les six macros qui
% suivent sont donc l'appareil critique, et non de la décoration.
%
% Les mêmes commandes sont reconnues par `scripts/render.mjs`, qui produit la
% vue de lecture du site. Une macro ajoutée ici doit l'être aussi là-bas,
% faute de quoi le rendu échouera bruyamment — ce qui est voulu : un rendu
% silencieusement faux est pire qu'un rendu absent.

\ProvidesPackage{grothendieck}[2026/08/08 Transcriptions du fonds Grothendieck]

\usepackage{amsmath,amssymb,amsthm}
\usepackage{mathrsfs}
\usepackage{stmaryrd}
% Les diagrammes commutatifs sont partout dans ce fonds : c'est la langue
% même de la géométrie algébrique de Grothendieck, et une transcription qui
% les rendrait en prose ne transcrirait rien.
\usepackage{tikz-cd}
% Français par défaut — c'est la langue des feuillets — mais l'anglais est
% chargé aussi : la traduction et le résumé sont des documents anglais, et la
% typographie française (espaces avant les deux-points) y serait une faute.
% Ces documents basculent par \selectlanguage{english} en tête de corps.
\usepackage[english,french]{babel}
\usepackage{fontspec}
\usepackage{geometry}
\geometry{a4paper,margin=25mm,marginparwidth=28mm}
\usepackage{marginnote}
\usepackage{xcolor}
% ulem, not soul. `soul`'s \st and \ul are built on a character-by-character
% scan that XeLaTeX's font machinery breaks: the document fails with an
% undefined control sequence rather than degrading. ulem does the same two jobs
% and is Unicode-engine safe. `normalem` stops it hijacking \emph, which the
% transcriptions use for Grothendieck's own emphasis.
\usepackage[normalem]{ulem}
\usepackage{hyperref}
\hypersetup{colorlinks=true,linkcolor=black,urlcolor=black,pdfborder={0 0 0}}

% --- Métadonnées du lot -----------------------------------------------------
% Reprises telles quelles par le script de rendu, qui les affiche en tête de
% la vue de lecture. Elles fixent ce que le fichier prétend couvrir : sans
% elles, une transcription flotte sans point d'attache dans un fonds de seize
% mille feuillets.

\newcommand{\folder}[1]{\gdef\grothFolder{#1}}
\newcommand{\batch}[1]{\gdef\grothBatch{#1}}
\newcommand{\leaves}[2]{\gdef\grothFirst{#1}\gdef\grothLast{#2}}
\newcommand{\foldertitle}[1]{\gdef\grothFoldertitle{#1}}
\gdef\grothFolder{?}\gdef\grothBatch{?}\gdef\grothFirst{?}\gdef\grothLast{?}\gdef\grothFoldertitle{}

% --- Le résumé --------------------------------------------------------------
%
% Il ouvre la lecture modernisée, sous le titre « Résumé », et il est composé
% en petit. Ce n'est pas un ornement : c'est le seul passage du document qui
% ne soit pas des mathématiques, et un lecteur doit pouvoir le reconnaître
% avant de l'avoir lu — puis le sauter s'il connaît déjà le sujet, ou s'y
% arrêter s'il ne le connaît pas. Le filet qui le suit marque la reprise du
% corps.

\newenvironment{resume}
  {\small\setlength{\parskip}{0.45em}}
  {\par\medskip\hrule\medskip}

% --- Mots-clés ---------------------------------------------------------------
%
% En anglais, en clôture du résumé de la lecture modernisée : le vocabulaire
% moderne sous lequel chercher ce que les feuillets construisent. Le manifeste
% du site (`npm run manifest`) les extrait pour étiqueter la cote entière —
% cette ligne est la source unique des tags, il n'y a pas de fichier à côté.
\newcommand{\keywords}[1]{\par\smallskip\noindent
  {\footnotesize\textsc{Keywords} --- \textit{#1}}\par}

% --- L'appareil critique ----------------------------------------------------

% Le numéro de feuillet, en marge. C'est l'ancre qui permet de régler un
% désaccord sur une formule en regardant *un* feuillet plutôt qu'une chemise
% de six cents. Le site s'en sert aussi pour tourner les pages du fac-similé
% quand on fait défiler la transcription.
\newcommand{\leaf}[1]{%
  \marginnote{\footnotesize\color{gray!70}\textsf{#1}}%
  \label{leaf:#1}\ignorespaces}

% Illisible : on ne devine pas. Un mot inventé qui se lit comme les autres est
% le pire résultat possible.
\newcommand{\ill}{\textcolor{red!60!black}{[\,\ldots\,]}}

% Lecture douteuse : le mot est donné, et signalé comme incertain.
\newcommand{\uncertain}[1]{\uline{#1}}

% Ajout de l'éditeur — une lettre manquante, un mot sous-entendu.
\newcommand{\add}[1]{\textcolor{blue!50!black}{[#1]}}

% Biffé par Grothendieck lui-même. Ce qu'il a rayé fait partie du document :
% une rature dit souvent où la pensée a changé de direction.
\newcommand{\struck}[1]{\sout{#1}}

% Note du transcripteur, sur l'état matériel du feuillet.
\newcommand{\note}[1]{\par\smallskip{\small\color{gray!60!black}\textit{[#1]}}\par\smallskip}

% Note marginale de Grothendieck, distincte du corps du texte.
\newcommand{\marginal}[1]{\par\smallskip\noindent\hspace{1em}%
  {\small\color{gray!60!black}#1}\par\smallskip}

% --- Titre ------------------------------------------------------------------

\AtBeginDocument{%
  \begin{center}
    {\large\bfseries Fonds Alexandre Grothendieck --- cote n\textsuperscript{o}~\grothFolder}\\[2pt]
    {\itshape\grothFoldertitle}\\[4pt]
    {\small feuillets \grothFirst--\grothLast\ (lot \grothBatch)}\\[2pt]
    {\footnotesize\color{gray!60!black}Transcription établie d'après le fac-similé numérisé
     par l'Université de Montpellier.\\
     Les lectures incertaines sont soulignées, les illisibles notées [\,\ldots\,],
     les ajouts entre crochets.}
  \end{center}
  \vspace{1em}}

\endinput


\folder{58}
\batch{4}
\pages{61}{66}
\dating{[à partir de 1964]}
\watermark{Édition de démonstration}
\foldertitle{Théorie de " Lefschetz-Grauert ". Applications en π₁, à Pic etc… : notes manuscrites (s.d.)}

\begin{document}

\page{61}
$C'_i$ :
\note{Encre bleue. Deux accolades à gauche regroupent les lignes : la première, marquée « $C'_i$ », embrasse les trois premières formules ; la seconde, marquée « $C_n$ », embrasse les deux dernières.}

$H^i(U, F(m)) = 0$ pour $m$ grand ($i > 0$ donné)

$\Updownarrow$

$\underline{H}^{i+1}_Z(F) = 0$ i.e.\ $H^{i+1}_z(F_z) = 0$ pour $z \in Z$

$\Downarrow$

$\mathrm{R}^i g_*(F) = 0$

$\Updownarrow$
\note{Cette flèche est surchargée de traits obliques ; on ne sait si elle est biffée ou seulement repassée.}

$C_n$ : $H^i(U, F(m)) = 0$ pour $m$ grand, $0 < i \leq n$ ($n \geq 1$ donné)

$\Updownarrow$

\struck{$H^i$} [\ill{} \uncertain{supposant} \uncertain{que} $F = g_*(F|U)$] $\operatorname{prof}_z F > n+1$ pour $z \in Z$
\note{La seconde double flèche est tracée plus épais, à l'encre plus sombre ; le « $H^i$ » biffé ouvre la ligne. La lecture « supposant que » est incertaine ; le crochet se ferme avant « $\operatorname{prof}_z$ ».}

Les conditions $C$ \uncertain{sont} les seules qui exigent quelque \uncertain{chose} sur \uncertain{la} \uncertain{profondeur} en les $z \in Z = X - U$ \ldots\ et non seulement aux pts de $U$ !
\note{Le texte s'arrête sur un trait ondulé tiré en travers de la page.}

\page{64}
$\underline{\mathcal{O}}_{\hat{C}}(1) \rightsquigarrow \underline{\mathcal{O}}_{\hat{C}_X}(1)$ : induit \uncertain{$0$} sur $C$ et sur $C'$
\note{Feuille de brouillon, encre bleue et noire, avec de petites figures (cercles pointés, un cercle traversé d'une courbe, une croix de deux branches) qu'on ne reproduit pas. L'indication « induit … sur $C$ et sur $C'$ » est écrite sous la formule, derrière une accolade ; son premier signe (« $0$ » ou « $\mathcal{O}$ ») est incertain.}

$\hat{C} =$ \uncertain{Spec} $S[\mathbf{z}]$
\note{Le mot lu « Spec » pourrait être autre chose (on attendrait « Proj » pour le cône projectif, cf.\ la page 66).}

$\underline{\mathcal{O}}_{\hat{C}}(1)$

$\operatorname{Pic}(\hat{C}) \to$ \quad noyau \ill{}

\begin{tikzcd}[column sep=small, row sep=small, nodes={font=\scriptsize}]
\operatorname{Pic}(X) \arrow[r, "\text{injectif}"'] \arrow[rr, bend left=20] & \operatorname{Pic}(\hat{C}') \arrow[r] \arrow[dr] & \operatorname{Pic}(C') \arrow[r] & \operatorname{Pic}(\widetilde{C'}) \\
 & & \text{noyau} &
\end{tikzcd}

\note{Sous $\operatorname{Pic}(X) \to \operatorname{Pic}(\hat{C}')$ : « [bijectif dans le cas $X$ loc.\ factoriel] » ; sous $\operatorname{Pic}(C') \to \operatorname{Pic}(\widetilde{C'})$ : « injectif [surjectif] dans le cas \ill{} factoriel] », le premier mot récrit. Sous « noyau », relié par la flèche oblique : « [engendré par l'image de $\mathcal{O}_X(1)$] » (« engendré par » incertain). Un pointillé descend aussi de $\operatorname{Pic}(C')$ vers « noyau ». À gauche, au-dessus de $\operatorname{Pic}(X)$, « Pic C » biffé, avec une flèche montante.}

$\operatorname{Pic}(X)/\mathbf{Z}\xi \subset \operatorname{Pic}(\hat{C}')/\mathbf{Z}\xi$

$\boxed{\operatorname{Pic}(C) = 0 \ ??}$

\begin{tikzcd}[column sep=small, row sep=small, nodes={font=\scriptsize}]
 & \mathbf{Z}\xi \arrow[r, no head] \arrow[d] & \mathbf{Z}\xi \arrow[d, "(?)"] & \\
0 \arrow[r] & \operatorname{Pic}(X) \arrow[r] & \operatorname{Pic}(\hat{C}') \arrow[d] & \\
 & \operatorname{Pic}(C) \arrow[r] & \operatorname{Pic}(C') \arrow[r] & \operatorname{Pic}(\widetilde{C'})
\end{tikzcd}

\note{Les deux $\mathbf{Z}\xi$ sont reliés par un signe « $=$ » ; un « $0$ » au bas du cadre précédent descend vers le second. Le point d'interrogation qui marque sa flèche verticale est encerclé. Une flèche vers la gauche aboutit à $\operatorname{Pic}(\hat{C}')$ depuis la droite, sans source écrite ; la ligne du bas commence par un signe illisible. De part et d'autre de la flèche $\operatorname{Pic}(\hat{C}') \to \operatorname{Pic}(C')$ : à gauche « injectif [surjectif] », à droite « \struck{injectif} [surjectif] ».}

$\boxed{\operatorname{Pic}(C) = 0 \ ?}$ \quad Pic
\note{Second cadre, à droite, relié par un trait à un signe effacé.}

$g/f$ \qquad $1/g$

$g^m = u f^n$

\struck{$g = \ill{}$} \qquad \struck{$g^m / f$}

$g = u f^k$

\uncertain{$g^m f^n = u$}
\marginal{\ill{} \ill{} \ill{} \ill{} \ill{} \ill{} \ill{}}
\note{Colonne de formules en marge droite, près de la figure d'un cercle pointé traversé d'un trait vertical ; la dernière ligne, d'une autre encre, est d'une lecture incertaine. La note marginale, sur trois lignes, n'a pas été lue.}

\page{65}
$X - a$
\note{Feuille de brouillon, encre bleue ; de petites figures (cercles pointés, croix de deux ou trois droites, deux branches se coupant en un point marqué, un carré pointé dans un cadre) accompagnent les formules et ne sont pas reproduites.}

\[
  0 \to H^0_a(X, \mathcal{O}_X^*) \to H^0(X, \mathcal{O}_X^*) \to H^0(U, \mathcal{O}_X^*) \to H^1_a(X, \mathcal{O}_X^*) \to H^1(X, \mathcal{O}_X^*) \to \cdots
\]
\[
  \cdots \to H^1(U, \mathcal{O}_X^*) \to H^2_a(X, \underline{\mathcal{O}}_X^*) \to H^2(X, \underline{\mathcal{O}}_X^*)
\]
\note{Sous $H^1(X, \mathcal{O}_X^*)$ : « $= 0$ », en colonne. Les exposants de $H^1_a$ et de $H^2_a$ sont récrits sur un premier chiffre. Les deux derniers termes sont soulignés chacun d'une flèche vers la gauche.}

$A^* \to \tilde{A}^*$

$F$ \quad $H^0_x(F)$

loc.\ factoriel $\Rightarrow$ loc.\ quasi-factoriel $\Rightarrow$ parafactoriel

factoriel $\Rightarrow$ 1-factoriel $\Rightarrow$ 0-factoriel,
\note{« loc. » est écrit au-dessus de « quasi-factoriel » ; une double flèche oblique descend de « factoriel » vers « 1-factoriel », et une petite double flèche de « 1-factoriel » vers « 0-factoriel ».}

$\tilde{A} \to \tilde{A}_0$ \qquad \struck{\ill{}}

$H^1(X, \mathcal{O}_X) = 0$

$L$ \quad \struck{$L_B$} \qquad $\tilde{X}$ au-dessus de $X$

\struck{$1 - \ill{}$}

\uncertain{$1 - h'$} $= 2 - h' - h''$

$J \subset A \subset B$ \qquad $A/J \hookrightarrow B/JB$

\[
  1 - g = \sum (1 - g_i) - \sum n_s , \qquad 1 - g = \nu - \sum g_i - \sum n_s , \qquad g = 1 + \sum g_i + \sum n_s - \nu
\]

\[
  1 - g = 1 - g' - \sum n_s , \qquad g = g' + \sum n_s
\]
\[
  1 - g = 2 - g' - g'' - \sum n_s , \qquad g = (g' + g'') + \sum n_s - 1
\]
\note{Les deux cas sont séparés par un trait oblique : à gauche une seule composante de genre $g'$, à droite deux composantes de genres $g'$, $g''$. Près de la dernière formule : « \struck{$g' + g''$} », « $g' = g'' = 0$ », « $\sum g_i$ ».}

$k[x, y, z]/xy$
\note{À côté, la figure de deux droites sécantes coupées par une troisième, et celle de deux branches courbes se coupant en un point.}

\page{66}
$X$ schéma projectif sur $k$, $\Gamma(X, \mathcal{O}_X) \simeq k$

$C$ le cône projetant affine, a un \uncertain{sommet}

$\hat{C}$ le cône projetant projectif

\begin{tikzcd}[column sep=small, row sep=small, nodes={font=\scriptsize}]
X \arrow[r, "i"] \arrow[d] & X' \arrow[r, "\pi"] \arrow[d, "f"] & X \\
a \arrow[r] & \hat{C} &
\end{tikzcd}

\note{Entre les deux flèches verticales, en petits caractères : « \uncertain{section} \ill{} II 8.7.5 ». La ligne du bas se prolonge par un trait après $\hat{C}$.}

\struck{$\operatorname{Pic}(X) \xrightarrow{\pi^*} \operatorname{Pic}(X')$, donc}
[\uncertain{car} $X'$ \ill{} \ill{} \uncertain{fibré} \ill{} sur $X$]
\struck{[\ill{}]}
\note{Le début de ligne est barré d'un trait ondulé ; l'explication entre crochets, derrière une accolade, se rapporte à la ligne suivante.}

donc $\operatorname{Pic}(X') \xrightarrow{\ \sim\ } \operatorname{Pic}(X) \times \mathbf{Z}$
\note{Au-dessus de la flèche, un indice illisible.}

$i^*\colon \operatorname{Pic}(X') \to \operatorname{Pic}(X)$ \uncertain{a} \uncertain{comme} noyau $\simeq \mathbf{Z}$, engendré par $\underline{\mathcal{O}}_{X'}(1)$

$\operatorname{Pic}(\hat{C}) \to \operatorname{Pic}(X')$ injectif [car $f_*(\mathcal{O}_{X'}) = \underline{\mathcal{O}}_{\hat{C}}$] par II 8.8.6.

Je \uncertain{m'aperçois} que
\[
  \boxed{\operatorname{Pic}(\hat{C}) \simeq \mathbf{Z} \xrightarrow{\ \sim\ } \operatorname{Ker}\bigl[\operatorname{Pic}(X') \xrightarrow{i^*} \operatorname{Pic}(X)\bigr]}
\]
\note{Dans le cadre, « $\simeq \mathbf{Z}$ » est récrit en surcharge, et un trait vertical le sépare de la flèche.}

Soit $C_X$ la \uncertain{partie} de $X'$ au-dessus de $C$, c'est un fibré \struck{\ill{}} \uncertain{en} \struck{\ill{}} 1 sur $X$.
\note{Le « 1 » est écrit en surcharge, à l'encre plus sombre, après un mot biffé.}

Je dis que :
\[
  \operatorname{Pic}(X) \xrightarrow{\ \sim\ } \operatorname{Pic}(C_X) \quad (?)
\]
\marginal{faut-il supposer $X$ localement factoriel ??}
\note{Note marginale gauche, oblique, reliée au texte par une flèche. Un astérisque surmonte « Pic » à gauche de la flèche.}

\uncertain{Cet} \ill{} \uncertain{est} \ill{} injectif, car \uncertain{son} \uncertain{composé} avec $j^*\colon \operatorname{Pic}(C_X) \to \operatorname{Pic}(X)$, \ill{} \ill{}, \struck{\ill{} \ill{} \ill{} de $C_X$ sur} l'identité. Pour la \emph{surjectivité}, on \uncertain{voit} que \ill{} \uncertain{au-dessus} d'un \uncertain{ouvert} \uncertain{convenable} de $X$,
\note{La phrase s'interrompt sur la virgule ; le reste de la page est blanc. C'est la dernière page du dossier.}

\end{document}
